\section{Background}

Decision-making problems in engineering, management, healthcare, supply-chain selection, and public-sector planning often involve multiple alternatives and several conflicting criteria. Multi-criteria decision-making (MCDM) methods provide systematic tools for comparing alternatives when no single criterion is sufficient. Among these methods, the Technique for Order Preference by Similarity to Ideal Solution (TOPSIS) is widely used because it is intuitive, mathematically simple, and easy to communicate to decision makers. TOPSIS assumes that the best alternative should be close to the positive ideal solution and far from the negative ideal solution.

In many practical settings, decision-maker judgments are uncertain. A decision maker may not be able to state that an alternative has an exact score of 7 on a criterion, but may be comfortable saying that the rating is approximately between 6 and 8, with 7 as the most likely value. Fuzzy TOPSIS handles such uncertainty by representing ratings and weights as fuzzy numbers. This thesis uses triangular fuzzy numbers of the form \((l,m,u)\), where \(l\) is the lower value, \(m\) is the most plausible or modal value, and \(u\) is the upper value.

Group decision-making adds another challenge. Instead of relying on a single decision maker, a group of decision makers evaluates the alternatives. Group decision-making is useful because it can combine expertise and reduce individual error. However, it also creates a vulnerability: if some decision makers are biased or coordinated, the final group ranking can be manipulated.

\section{Motivation}

Classical fuzzy TOPSIS assumes that decision-maker inputs are sincere. In real decision environments, this assumption may be too strong. Decision makers can be influenced by personal preference, organizational pressure, conflict of interest, incomplete knowledge, or deliberate manipulation. If several decision makers coordinate their ratings, they can promote a weak target alternative by assigning high ratings to it and low ratings to competing alternatives.

This issue is important because fuzzy TOPSIS is often used in high-impact decision contexts. Supplier selection, healthcare resource allocation, institutional ranking, technology selection, and policy evaluation all depend on the reliability of group judgments. A method that performs well under clean data may not be trustworthy under biased group input.

The motivation of this research is therefore to make fuzzy TOPSIS more resistant to biased and coordinated decision-maker behavior. Instead of assuming every decision maker is equally reliable, this thesis estimates reliability from rating behavior and reduces the influence of suspicious patterns.

\section{Problem Statement}

The main problem addressed in this thesis is targeted rank manipulation in fuzzy TOPSIS group decision-making. The attack setting is:

\begin{enumerate}
    \item A clean group of decision makers evaluates a set of alternatives.
    \item The clean fuzzy TOPSIS ranking is computed.
    \item The lowest-ranked alternative is selected as the attack target.
    \item A fraction of decision makers becomes contaminated.
    \item Contaminated decision makers assign artificially high ratings to the target and low ratings to other alternatives.
    \item The method is evaluated by whether it preserves the clean target rank.
\end{enumerate}

The core research question is:

\begin{quote}
Can reliability-aware fuzzy TOPSIS methods preserve the original group decision when a subset of decision makers provides coordinated biased ratings?
\end{quote}

\section{Research Questions}

This thesis investigates the following research questions:

\begin{enumerate}
    \item How vulnerable is classical fuzzy TOPSIS to coordinated decision-maker contamination?
    \item Does bootstrap bagging alone improve robustness?
    \item Does reliability-weighted sampling improve over plain bagging?
    \item Does reliability-weighted fuzzy aggregation suppress biased decision-maker influence more effectively?
    \item Can entropy, variance consistency, and clone/agreement-pattern signals improve robustness under high structured contamination?
    \item How do the proposed methods compare with robust aggregation, consensus filtering, individual-rank aggregation, and group-ideal TOPSIS baselines?
\end{enumerate}

\section{Objectives}

The objectives of this thesis are:

\begin{enumerate}
    \item To implement a classical fuzzy TOPSIS baseline for group decision-making.
    \item To correct and evaluate a true bootstrap bagging version of fuzzy TOPSIS.
    \item To develop reliability-aware fuzzy TOPSIS methods that reduce the influence of suspicious decision makers.
    \item To propose three main methods: M4, M6, and M7/\methodname.
    \item To evaluate the methods on synthetic and real/pseudo-real datasets.
    \item To compare the proposed methods against external comparator baselines.
    \item To provide a reproducible evidence package and a public API scaffold for future users.
\end{enumerate}

\section{Contributions}

The main contributions are:

\begin{itemize}
    \item A structured vulnerability analysis showing that classical fuzzy TOPSIS can be manipulated by coordinated decision-maker bias.
    \item A corrected bootstrap fuzzy TOPSIS baseline that uses true sampling with replacement.
    \item Reliability-weighted bagging (M4), where reliable decision makers are sampled more frequently.
    \item Reliability-weighted aggregation (M6), where decision-maker ratings are directly weighted by reliability during fuzzy aggregation.
    \item \methodname{} (M7), an entropy-aware reliability-weighted robust fuzzy TOPSIS framework using entropy, variance consistency, and clone/agreement-pattern signals.
    \item A full attack-fraction evaluation showing method behavior from low to high contamination.
    \item External comparator evaluation including median TOPSIS, trimmed-mean TOPSIS, MAD-consensus TOPSIS, individual-Borda TOPSIS, and a Huang--Li group-ideal TOPSIS adaptation.
    \item A public web/API scaffold for uploading datasets and running M4, M6, and M7.
\end{itemize}

\section{Scope and Limitations}

The thesis focuses on structured and statistically distinguishable coordinated bias. It does not claim to detect all forms of human bias. Honest disagreement among decision makers is normal and should not automatically be penalized. The limitation appears when malicious decision makers deliberately mimic honest variation. If biased decision makers are statistically indistinguishable from honest decision makers, unsupervised reliability estimation cannot guarantee detection without additional information such as historical trust, external validation, or supervised labels.

\section{Organization of the Thesis}

Chapter 2 reviews TOPSIS, fuzzy TOPSIS, group decision-making, consensus models, expert weighting, and the research gap. Chapter 3 presents the methodology and all seven methods. Chapter 4 describes datasets, attack design, metrics, repeated runs, and comparator baselines. Chapter 5 presents the experimental results. Chapter 6 discusses interpretation, limitations, and practical implications. Chapter 7 concludes the thesis and outlines future work. The appendix describes the planned public web and API platform.
